% Generated draft: requires mathematical and source review.
\documentclass[11pt]{article}
\usepackage[margin=25mm]{geometry}
\usepackage{fontspec}
\setmainfont{Latin Modern Roman}
\usepackage{amsmath,amssymb,mathtools,unicode-math}
\setmathfont{Latin Modern Math}
\usepackage{enumitem,needspace,xurl,hyperref}
\setlength{\parindent}{0pt}
\setlength{\parskip}{5pt}
\setlength{\emergencystretch}{3em}
\newcommand{\N}{\mathbb{N}}
\newcommand{\Z}{\mathbb{Z}}
\newcommand{\Q}{\mathbb{Q}}
\newcommand{\R}{\mathbb{R}}
\newcommand{\C}{\mathbb{C}}
\newcommand{\F}{\mathbb{F}}
\newcommand{\A}{\mathbb{A}}
\newcommand{\PP}{\mathbb{P}}
\newcommand{\E}{\mathbb{E}}
\newcommand{\eps}{\varepsilon}
\newcommand{\dd}{\,\mathrm{d}}
\begin{document}
\section*{UnsolvedMath ID: 30002795}
\textbf{Source problem code:} catalogue entry
\subsection{Definitions and mathematical statement}
Let \(k\) be a field of characteristic zero and let \(r\geq 4\) be an integer. Give \(S=k[x_1,\ldots,x_r]\) its standard grading, with \(\deg x_j=1\). Under the usual graded Artinian interpretation of the supplied question, consider
\[
A=S/(f_1,\ldots,f_r)=\bigoplus_{d\geq 0}A_d,
\]
where the \(f_j\) are homogeneous polynomials of positive degree forming a regular sequence, and \(A\) is finite-dimensional over \(k\). Here regular sequence means that the generated ideal is proper and, for each \(j\), the image of \(f_j\) is a non-zero-divisor in \(S/(f_1,\ldots,f_{j-1})\). Such an \(A\) is a standard graded Artinian complete intersection.

The algebra \(A\) has the \emph{weak Lefschetz property} if there exists a single element \(\ell\in A_1\) such that, for every integer \(d\geq 0\), multiplication induces a map
\[
\times\ell:A_d\longrightarrow A_{d+1}
\]
of maximal rank, meaning
\[
\operatorname{rank}_k(\times\ell)=\min\{\dim_k A_d,\dim_k A_{d+1}\}.
\]
\textbf{Question.} For every field \(k\) of characteristic zero, every integer \(r\geq4\), and every homogeneous regular sequence as above, does \(A\) have the weak Lefschetz property?

\textit{Qualification of the formulation.} The supplied short statements omit characteristic, homogeneity, and Artinian hypotheses. The supplied literature assessment explicitly identifies the intended question as concerning characteristic-zero Artinian complete intersections; homogeneity and the standard grading are made explicit here to formulate the usual graded weak Lefschetz property. The source record does not explicitly specify the precise field scope or grading hypotheses, so these aspects require human confirmation against the report. The original extraction also asked what happens for more than \(r+1\) forms. The supplied correction identifies that as a distinct question and removes it; it is not included in the quantified question above.
\subsection{Short English statement}
Does every standard graded Artinian complete intersection in at least four variables over a characteristic-zero field admit one linear form whose multiplication maps between all consecutive graded components have maximal rank? This formulation follows the supplied characteristic-zero Artinian qualification and makes the usual homogeneous grading explicit, subject to source review. The supplied open status is dated curation, not an independently verified conclusion.
\subsection{Sources}
\begin{itemize}
\item Local registry: \texttt{data/unsolvedmath/problems.json}, permanent UnsolvedMath ID \(30002795\), source problem code \texttt{catalogue entry}, title \emph{Weak Lefschetz Properties of Complete Intersections}, in the Oberwolfach Reports Open Problems collection. The record gives proposed year 2015 and names no proposer.
\item Supplied report citation: \emph{Mini-Workshop: Ideals of Linear Subspaces, Their Symbolic Powers and Waring Problems} (2015), Oberwolfach Reports, pp.\ 489--532, DOI \texttt{10.4171/owr/2015/9}. \url{https://doi.org/10.4171/owr/2015/9}. The supplied background describes these reports as documenting mathematical workshops and their open-problem discussions.
\item Supplied literature source: \emph{On the weak Lefschetz property for height four complete intersections}. \url{https://arxiv.org/abs/2212.09890}.
\item Literature assessment checked on \texttt{2026-08-21}: the supplied curation labels the characteristic-zero question for arbitrary Artinian complete intersections in at least four variables as open and reports later work on restricted height-four and equigenerated cases. This is supplied dated curation, not independently verified; its status and cited claims require independent checking.
\item The record labels its statement \texttt{corrected\_verified} and reports that the official report presents the complete-intersection question separately. Its verification note says the extraction appended a distinct question about more than \(r+1\) generators, whose unrestricted form already has negative examples, and that this clause was therefore removed. These are supplied source-verification claims, not independent verification performed for this draft.
\end{itemize}
\end{document}
